% Part: applied-modal-logics
% Chapter: temporal-logic
% Section: properties-accessibility

\documentclass[../../../include/open-logic-section]{subfiles}

\begin{document}

\olfileid{aml}{tl}{acc}

\olsection{కాలిక చట్రాల ధర్మాలు}

మన కాలిక నమూనాలు సంబంధం~$\prec$పై ఎలాంటి షరతులూ
విధించనందున, పరిచయమైన స్వీకృతాల్లో అన్ని నమూనాల్లో
నిలిచేది~$K$ మాత్రమే; లేదా దాని సదృశ రూపాలైన
$K_{\Gtemp}$,~$K_{\Htemp}$ ఇవి:
\begin{align*}
\tag{$K_{\Gtemp}$}  \Gtemp (p \to q) & \to (\Gtemp p \to \Gtemp q)\\
\tag{$K_{\Htemp}$}  \Htemp (p \to q) & \to (\Htemp p \to \Htemp q)
\end{align*}

అయితే నమూనాల్లో కాలాన్ని సూచించే విధానంపై కఠినమైన
షరతులు, ఉదాహరణకు $\prec$ రేఖీయ క్రమం కావాలనే షరతు,
విధిస్తే ఇతర చెల్లుబాటులూ వస్తాయి.

\begin{table}[t]
  \begin{tabular}{| p{.48\textwidth} || p{.48\textwidth} |}
    \hline
    {\emph{$\prec$ ఇలా ఉంటే \dots}} & {\emph{$\mModel{M}$లో ఇది సత్యం \dots:}} \\
    \hline \hline
    \emph{సంక్రామకం}: \newline
    $\forall u \forall v \forall w ((u \prec v \land v \prec w) \lif u \prec w)$ &
    $\Ftemp \Ftemp p \lif \Ftemp p$  \\
    \hline
    \emph{రేఖీయం}: \newline
    $\forall w \forall v (w \prec v \lor w = v \lor v \prec w)$ &
    $(\Ftemp \Ptemp  p \lor \Ptemp \Ftemp p) \lif (\Ptemp p \lor p \lor \Ftemp p)$\\
    \hline
    \emph{సాంద్రం}: \newline
    $\forall w \forall v (w \prec v \to \exists u(w \prec u \land u \prec v))$ &
    $\Ftemp p \lif \Ftemp \Ftemp p$ \\
    \hline
    \emph{గతంలో అంచులేనిది}: \newline
    $\forall w \exists v( v \prec w)$ &
    $\Htemp p \to \Ptemp p$ \\
    \hline
    \emph{భవిష్యత్తులో అంచులేనిది}: \newline
    $\forall w \exists v( w \prec v)$ &
    $\Gtemp p \to \Ftemp p$ \\
    \hline
  \end{tabular}
  \caption{కొన్ని కాలిక చట్ర అనురూప ధర్మాలు.}
  \ollabel{tab:correspondence}
\end{table}

\olref{tab:correspondence}లోని అనేక ధర్మాలు కాలాన్ని
సూచించే నమూనాకు కావాల్సినవిగా అనిపించవచ్చు. అయితే
$\prec$ సంబంధంపై మనకు నచ్చిన షరతులు విధించగలిగినా,
ప్రతి షరతుకూ కాలిక తర్క భాషలో వ్యక్తం చేయగల
\tetoken{సూత్రం}{!!{formula}s} అనురూపంగా ఉండదని
గమనించాలి. ఉదాహరణకు, అస్వావర్తనత్వం, అంటే
$\forall w \lnot (w \prec w)$, కోసం కాలిక తర్కంలో
అనురూప సూత్రం లేదు.

\end{document}
